\documentclass{article}
\begin{document}

\title{Optimal Control of an Epidemic Spreading on a Network}
\author{Anonymous Author}

\section{Introduction}
We study the spread of an infectious disease over a contact network.
Previous work mostly assumed homogeneous mixing, which fails to capture
real population structure. Our model extends the classic SIR framework.

\section{Problem Statement}
Given a graph $G=(V,E)$ and initial infections, find a control policy
that minimizes total infections subject to a budget constraint.

\section{Model Formulation}
The dynamics follow a degree-based mean-field approximation:
\[
\frac{dS_k}{dt} = -\beta k S_k \Theta(t), \qquad
\frac{dI_k}{dt} = \beta k S_k \Theta(t) - \gamma I_k.
\]
Here $\Theta(t)$ is the probability a random edge points to an infected node,
defined by
\[
\Theta(t) = \frac{1}{\langle k\rangle}\sum_{k'} k' P(k') \frac{I_{k'}}{N}.
\]

\section{Methodology}
We discretize time and solve the resulting ODE system with an
implicit Euler scheme. Parameters: $\beta=0.4$, $\gamma=0.1$,
$\langle k\rangle=6$.

\section{Numerical Experiments}
We simulate a network of $N=10^4$ nodes. The controlled policy reduces
peak infected fraction by a noticeable margin compared with no control.
Results are summarized in Figure~\ref{fig:curve} and Table~\ref{tab:res}.

\begin{figure}[htbp]
  \centering
  \caption{Infected fraction over time under control vs. no control.}
  \label{fig:curve}
\end{figure}

\begin{table}[htbp]
  \centering
  \caption{Peak infected fraction under different budgets.}
  \label{tab:res}
  \begin{tabular}{lc}
    \toprule
    Budget & Peak Infected \\
    \midrule
    0      & 0.42 \\
    50     & 0.31 \\
    100    & 0.27 \\
    \bottomrule
  \end{tabular}
\end{table}

\section{Conclusion}
The degree-based mean-field model captures heterogeneity and enables
efficient control optimization. Future work should validate against
real mobility data.

\end{document}